\documentclass[11pt, letterpaper]{article}
\usepackage[margin=.75in]{geometry}
\usepackage{setup}

% ==========================================
% Define parameters for this specific recital
% ==========================================
\setperformers{\emph{Junrui Liu}}
\setdatetimeplace{Sunday, May 8, 2026}{2:00 p.m.}{Storke Tower}

\begin{document}
\begin{center}
    {\departmenthead} \\[1em]
    {\Large \recitaltype} \\[.5em]
    {\Large \textbf{\theperformers}} \\[.5em]
    {\small \thedatetimeplace} \\[1em]
\end{center}

\ruleline{program}

\vspace{1em}

\newcommand{\motif}[1]{\item[] {\centering \textit{#1} \par\vspace{-.5em}}}

\begin{programlist}
    \motif{To Flow}
    \item
    \piece{Prelude (from \emph{Fugue and Allegro in E-flat major, BWV 998})}{J. S. Bach (1685--1750)}
    \piece{Image 1: Brouillard (2016)}{Stefano Colletti (b. 1973)}


    \motif{To Remember}
    \item
    \piece{Sicilienne Ronde, Op. 22 (2005)}{John Gowens (b. 1957)}
    \piece{Reflection}{Robert Byrnes (1949-2004)}


    \motif{To Sing}
    \item
    \piece{Chartres (1969)}{Roy Hamlin Johnson (1929-2020)}
    \piece{Musetta's Waltz (from \emph{La Bohème})}{Giacomo Puccini (1858-1924)}
    \arranger{Sally Slade Warner}

    \motif{To Dream}
    \item
    \piece{Sonatine (2001)}{Stefano Colletti}
    \piece{Modal Nocturne (1996)}{Geert D'hollander (b. 1965)}
\end{programlist}

\vspace{2em}

\begin{center}
    \small \it Each section will be announced by the striking of one or more bells
\end{center}

\hrulefill

\textbf{Welcome to Storke Tower and its Carillon!}

{\small
    The 61-bell UCSB carillon was dedicated in 1969. It was a gift of Thomas Storke, then publisher of the \textit{Santa Barbara News-Press}. The instrument, with bells cast by the Dutch bell foundry Petit \& Fritsen, is one of the finest carillons in the world with regard to bell tuning and sound quality, number of bells, and location (not near a busy street and near the ocean). The bells range in weight from about 13 pounds to 2.5 tons.
    Our instrument is one of six carillons in California, with the others being at UC Berkeley, UC Riverside, Stanford University, Christ Cathedral (formerly the Crystal Cathedral) in Garden Grove, and Trinity Cathedral in San Jose. There are over 650 carillons in the world and over 180 in North America. The carillon art is a growing one, for new instruments are installed every year. \par
}

\textbf{About the Artist}

\begin{small}
    \textbf{Junrui Liu} is a Computer Science PhD candidate. His interest in the carillon was sparked by a series of live-streamed recitals by University Carillonist Wesley Arai during the pandemic. Captivated by the instrument’s unique sound, he joined Arai's studio as soon as a vacancy arose. Junrui's musical background also includes organ (if he chances upon one to play) and piano. This is Junrui's third year studying the carillon. \par
\end{small}

\vfill
\calloutbox{
    \begin{center}
        \small
        \textit{Please join us for our upcoming carillon recitals:} \\[.5em]
        Sunday, May 17, 2026 at 2:00pm -- Wesley Arai, University Carillonist \\
        Sunday, May 31, 2026 at 10:00am -- UCSB Carillon Studio Student Recital \\
    \end{center}
}

\vfill

\newpage

\textbf{Program Notes}
% \newcommand{\divider}{\begin{center}
%         \scalebox{0.7}{\adforn{21}} % Reduces total size to 70%
%     \end{center}}
\newcommand{\divider}{\vspace{.5em}}

\begin{small}
    \piecetitle{Prelude, Fugue and Allegro in E-flat major, BWV 998} was likely written for the Baroque lute, a precursor to the modern guitar. The Prelude movement is similar in style to many of the preludes in Bach's \emph{Well-Tempered Clavier}. It primarily features \emph{arpeggios}---Italian for ``playing the harp''---which arrange chordal notes in succession rather than simultaneously, giving the piece a flowing feel. The music modulates through several keys fifths before building toward a climax with a brilliant passage that transports from E-flat major to A-flat minor, two distantly related keys. Finally, it brings back the original theme in the home key over a pedal point and closes with a serene suspension.\\[.2em]
    %
    \piecetitle{Brouillard} {\it (``Fog'')} by French composer and carillonneur Stefano Colletti captures the essence of French Impressionism. The composition begins with swirling motions and ambiguous harmonies. The dense fog gradually gives way to a shimmering middle section, dotted with broad-brush strokes of richly colored chords in the lower voice and delicate \emph{glissandi} in the treble, like faint rays of sun piercing through. In order to maintain a light, glassy timbre, these glissandi (derived from the French word \emph{glisser} meaning ``to slide'') are performed with a pianistic finger technique, in contrast to the more common way of using fists to strike the carillon's baton keys. The work then returns to the spiraling motion of the opening, as if the ocean wind blows the mist back again.

    \divider

    \piecetitle{Sicilienne Ronde, Op. 22} (2005) is written by American carillonneur and organist John Gouwens as an homage to Ronald Barnes (1927–1997). Barnes was an eminent composer and carillonist whose work helped establish the American school of carillon composition and performance. One can discern traces of Barnes' style in this rondo-form piece through its harmonic language (e.g., frequent juxtaposition of themes in distantly related keys) and the song-like melodies used throughout. The final appearance of the ``A'' theme employs a compositional device common to Barnes' works, where the lyrical melody is placed in the pedals beneath a texture of rippling arpeggios in the manual.\\[.2em]
    %
    \piecetitle{Reflection} by American carillonneur Robert Byrnes was dedicated to the memory of Frank Péchin Law, also a carillonneur. The ternary-form composition unfolds gradually: first, sparse chords and light arpeggios establish the recurring harmonic backbone of the opening section, and introduce a bittersweet melody that seems to yearn and reach for something just out of grasp. The underlying harmony may sound familiar, for it is also known as the ``'50s progression'' found in many popular songs of the era. The serenity of the opening section is soon disturbed by a ``turning out the lights'' moment, when the secondary theme is replayed in the parallel minor key, darkening the tonal color. The ominous clouds soon brew into a roaring, unsettling spiral of non-functional harmony; here, the manual anxiously arpeggiate at twice the speed of the opening. The thunderstorm eventually subsides, washing back into the hauntingly beautiful theme of the first section.

    \divider

    \piecetitle{Chartres} (1969), composed by Roy Hamlin Johnson, a prolific American composer for the carillon, is adapted from a 15th-century French hymn tune (nowadays also associated with the hymn \emph{Saw You Never, in the Twilight}). It is in ``Part Three: Epiphany'' of his \emph{A Carillon Book for the Liturgical Year}, a collection of compositions meant to be played throughout the Christian year. The music begins with the tune set against a sparse, descending accompaniment, which then leads into \emph{canons} (voices played in imitations of one another). Increasingly, the music becomes agitated, crescendoing into a thunderous \emph{fortississimo} (fff) before fading away in staggered voices. The original tune reappears at last, and dies away with a haunting, mysterious amen.\\[.2em]
    %
    \piecetitle{Musetta's Waltz}, also known by its original aria title \enquote{Quando m'en vo'}, is a soprano staple from Puccini's 1896 opera \emph{La bohème}. It is sung by the spirited Musetta, who hopes to reclaim the attention of her on-again, off-again boyfriend, Marcello, by making him jealous. ``When walking alone on the streets / People stop and stare / And examine my beauty / From head to toe / And then I savor the cravings…'' as the lyrics go.

    \divider

    \piecetitle{Sonatine} (2001) is noted for its use of a carillon playing technique called \emph{tremolando}, where two bells (usually separated by a wide interval) are struck in a rapidly alternating fashion. This creates the illusion that a note is sustained for a longer duration, making it possible to ``sing'' a slow, lyrical melody as a continuous stream on an otherwise percussive instrument. Thematically, the central section evokes the works of French Impressionist composer Maurice Ravel, and might especially remind us of the ``Adagio Assai'' movement of Ravel’s \emph{Piano Concerto in G Major}.\\[.2em]
    %
    \piecetitle{Modal Nocturne} (1996) begins with the manual playing an airy accompaniment that oscillates between two broken chords, G major and d minor, which place this section in the G Mixolydian mode and  explains the ``modal'' in the title. A song-like melody soon joins on the lower bells. After a tonally wandering transition, the main theme reappears and then recast in the B Dorian mode. The middle section is a sweet and tranquil minuet. The final section sees the return of the main theme, this time voiced by the delicate upper bells and accompanied by a measured glissando in the inner voice. Together, they slowly build momentum, eventually cresting with an arresting double suspension before everything fades away into a tender dream.

\end{small}

\vfill


\end{document}